\subsection{LEAST UPPER BOUND OF A SET OF UNIFORMITIES}

Every family \((\mathcal{U}_t)_{i\in I}\) of uniformities on a set \(X\) has a least upper bound \(\mathcal{U}\) in the ordered set of all uniformities on \(X\) ; we have only to apply Proposition 4 of no. 3, taking \(Y_t\) to be the set \(X\) with the uniformity \(\mathcal{U}_t\) , and \(f_t\) to be the identity mapping \(X\to Y_t\) . The topology induced by \(\mathcal{U}\) is just the least upper bound of the topologies induced by the \(\mathcal{U}_t\) .

It follows also from Proposition 4 of no. 3 that if X is not empty and if \(\mathcal{U}_{t}\) is the filter of entourages of \(\mathcal{U}_{t}\) , then the filter of entourages of \(\mathcal{U}\) is the least upper bound of the filters \(\mathcal{U}_{t}\) (Chapter I, § 6, no. 2).

Example. If \(\varpi\) is any finite partition \((\mathbf{A}_i)_{1 \leqslant i \leqslant n}\) of a non-empty set \(\mathbf{X}\) , the set \(\mathbf{V}_{\overline{\omega}} = \bigcup_{i} (\mathbf{A}_i \times \mathbf{A}_i)\) by itself constitutes a fundamental system of

entourages of a uniformity \(\mathfrak{U}_{\varpi}\) on X (§ 1, no. 1, Example 2); the uniformity of finite partitions on X (no. 2, Remark 1) is then the least upper bound of the uniformities \(\mathfrak{U}_{\varpi}\) .

Remark. A family \((\mathcal{U}_{i})\) of uniformities on X also has a greatest lower bound in the ordered set of all uniformities on X, namely the least upper bound of the set of all uniformities on X which are coarser than each of the \(U_{i}\) (such uniformities exist, since the set of all uniformities on X has a least element). But (supposing X not empty) the filter of entourages of this uniformity is not necessarily the intersection of the filters of entourages of the \(U_{i}\) , because this latter filter need not satisfy axiom \((\mathrm{U}_{\mathrm{III}})\) (Exercise 4).

